Every subsequent bound is now evaluated on this identical path, including
the centered Hilbert--Schmidt, angle and signed finite-residual controls.
For compact notation put $n=2q+1$, $a_q=2q-1$, and define
\[
\begin{split}
 T_q(x)&=\tfrac12d_x\|U_x-W_x\|_{1,M_x},\\
 S_q(x)&=\sqrt{\left(\operatorname{Tr}(C_x^2)
       -\frac{(\operatorname{Tr}C_x)^2}{n}\right)
                    \operatorname{Tr}((U_x-W_x)^2)},\\
 z_x&=a_q\sqrt{1-e^{-F(x)/a_q}},\qquad A_q(x)=z_xd_x,\\
 J_q&=\int_0^1\min\{S_{\rm AW}(x),T_q(x),S_{\rm SP}(x),
                                      S_q(x),A_q(x)\}\,dx,\\
 I_q&=\int_0^1\min\{d_x,S_{\rm AW}(x)/z_x,T_q(x)/z_x,
                         S_{\rm SP}(x)/z_x,S_q(x)/z_x\}\,dx.
\end{split}
\tag{ACM11a}
\]
These symbols are exactly $E_q,T_q,O_q,S_q,A_q,J_q,I_q$ of
(RMT9)--(RMT13), with $E_q=S_{\rm AW}$ and $O_q=S_{\rm SP}$.
Equations (ACM1)--(ACM5) prove the operator and measure identities
needed for that substitution. In particular they retain the actual
inverse metric in the signed density $v(U_x-W_x)M_x^{-1}v^*$.

Here is the proof of the additional terms at this earlier use site.
Trace zero in (ACM10) allows subtraction of $\operatorname{Tr}C_x/n$
from $C_x$; Hilbert--Schmidt Cauchy--Schwarz gives $|F'|\leq S_q$.
The midpoint argument already proved gives $|F'|\leq T_q\leq S_{\rm SP}$.
Retain both complete $q$-entry angle lists of the minimum-section
pairs $(q-1,2q)$ and $(q,2q-1)$. The second pair has a designated
zero angle by the original ambient dimension. Each additional zero
remains in its list. The exact projection-difference spectra give
$|F'|\leq d_x\sum\sin\theta$ over the two lists. Their cosine-square
determinants give $F=\sum-\log\cos^2\theta$; applying concavity of
$t\mapsto\sqrt{1-e^{-t}}$ to the remaining $2q-1$ entries gives
$|F'|\leq A_q$, exactly as proved with every multiplicity in
(RMT11). The designated zero still contributes its displayed zero
to both scalar sums.

The denominator in (ACM11a) is strictly positive on this actual path.
Indeed equality $F=0$ makes both crossed minimum sections equal:
their Gram difference is the squared norm of their difference,
and every relative eigenvalue is at least one. Equality of their
determinants makes every such eigenvalue one. In particular the
constant remainder $1$ would be orthogonal to $\chi\mathcal P_q$.
The polynomial $\chi^{\#_k}(S)=\overline{\chi(k-\overline S)}$
has degree $q$ and equals $\overline{\chi(S)}$ on $S=k/2+iy$.
That orthogonality would imply $\int|\chi|^2d\mu_x=0$, contradicting
the positive Gram of the nonzero polynomial $\chi$. This proves
$F(x)>0$ without inserting a new hypothesis. Compactness gives a
positive minimum. All remaining functions are continuous except
the intersection dimension, whose rank strata are Borel sets
defined by matrix minors. Thus every displayed integral is finite.
Integrating the five bounds on the same signed integrand proves
\[
 \boxed{|\Delta_{h,k}|\leq J_q
 \leq\int_0^1\min\{S_{\rm AW}(x),S_{\rm SP}(x)\}\,dx.}
 \tag{ACM11}
\]
The earlier minimum remains an explicit majorant of the current bound.
If $b_1\leq\cdots\leq b_n$ are the positive eigenvalues of
$D=M_0^{-1}M_1$ in $M_0$, direct inversion of
$M_x=M_0((1-x)I+xD)$ gives
\[
 c_j(x)=\frac{b_j-1}{1-x+xb_j},\qquad
 \int_0^1c_j(x)\,dx=\log b_j.
 \tag{ACM12}
\]
The order is preserved because the derivative with respect to $b_j$
is $(1-x+xb_j)^{-2}>0$. Termwise integration of the finite sum gives
\[
\begin{split}
 |\Delta_{h,k}|\leq J_q
 &\leq\min\left\{\mathcal L_q(D),\int_0^1T_q\,dx,
 \int_0^1S_{\rm SP}\,dx,\int_0^1S_q\,dx,\int_0^1A_q\,dx\right\},\\
 \mathcal L_q(D)&=\log\frac{b_{q+2}}{b_q}
       +2\sum_{j=1}^{q-1}\log\frac{b_{2q+2-j}}{b_j}
       \leq(2q-1)\log\frac{b_n}{b_1}.
\end{split}
\tag{ACM13}
\]
The $q=1$ sum is empty. The last inequality follows by bounding
each of its $1+2(q-1)$ logarithmic factors by $\log(b_n/b_1)$.

The nonlinear coordinate also receives every one of these controls.
For $B>0$ put
\[
 \mathscr H_{a_q}(B)=2\operatorname{arcosh}(e^{B/(2a_q)}),\quad
 \Psi_{a_q}(t)=2a_q\log\cosh(t/2)\quad(t\geq0).
 \tag{ACM14}
\]
Their two inverse composites on the positive half-lines are identities.
Direct differentiation gives
\[
 \mathscr H_{a_q}'(B)=\frac{1}{a_q\sqrt{1-e^{-B/a_q}}}.
\]
Dividing the five proved derivative bounds
by $z_x$ and integrating therefore gives
$|\mathscr H_{a_q}(F(1))-\mathscr H_{a_q}(F(0))|\leq I_q$.
Moreover $I_q\leq\int d_x=\log(b_n/b_1)$.

For every integer $p\geq0$ retain the following specified finite
signed construction, on the same source operators:
\[
\begin{split}
 B_x^\circ&=(1-x)I+xD,\qquad
 \alpha_x=\tfrac12(2(1-x)+x(b_1+b_n)),\qquad
 H_x^\circ=I-B_x^\circ/\alpha_x,\\
 C_x^{[p]}&=\alpha_x^{-1}\sum_{r=0}^p(H_x^\circ)^r(D-I),\\
 R_x^{[p]}&=(H_x^\circ)^{p+1}C_x,
 \qquad E_x^{[p]}=R_x^{[p]}-\operatorname{Tr}(R_x^{[p]})I/n,\\
 j_p&=\int_0^1\operatorname{Tr}(C_x^{[p]}(U_x-W_x))\,dx,\\
 e_p&=\int_0^1\sqrt{\operatorname{Tr}((E_x^{[p]})^2)
                         \operatorname{Tr}((U_x-W_x)^2)}\,dx.
\end{split}
\tag{ACM15}
\]
The finite geometric identity proves $C_x-C_x^{[p]}=R_x^{[p]}$.
Every factor is a real function of $D$, hence self-adjoint in $M_x$.
Trace zero removes the centered scalar, so Hilbert--Schmidt
Cauchy--Schwarz proves $j_p-e_p\leq\Delta_{h,k}\leq j_p+e_p$.
Writing $\vartheta=(b_n-b_1)/(b_n+b_1)<1$ and
$K_D=\sum_{j=1}^n(|b_j-1|/\min\{1,b_j\})^2$, the full residual
eigenvalues give
$\operatorname{Tr}((E_x^{[p]})^2)\leq\vartheta^{2p+2}K_D$.
Both $U_x,W_x$ dominate their range projections, whose intersection
has dimension at least one, so $\operatorname{Tr}(U_xW_x)\geq1$.
Equation (ACM10) then gives
$e_p\leq\vartheta^{p+1}\sqrt{K_D(8q-6)}\to0$.
This retains the exact signed center and its controlled residual,
with no assumption of uniformity in the tensor order.

Put $F_0=F(0)$ and $H_0=\mathscr H_{a_q}(F_0)$. For every $p\geq0$
define the actual simultaneous endpoint functions
\[
\begin{split}
 \beta_p^-&=\max\{0,F_0-J_q,
       \Psi_{a_q}(\max\{0,H_0-I_q\}),F_0+j_p-e_p\},\\
 \beta_p^+&=\min\{F_0+J_q,
       \Psi_{a_q}(H_0+I_q),F_0+j_p+e_p\}.
\end{split}
\tag{ACM16}
\]
The preceding additive, nonlinear and signed inequalities prove
$\beta_p^-\leq F(1)\leq\beta_p^+$ by taking their intersection.
Both original angle lists, source masses, relation maps and arithmetic
Taylor unit remain fixed. At $D=cI$ all diameter and residual terms
vanish: each quotient volume is multiplied by the same scalar power
$q$, which cancels with the four signs. The formulas then give
$F(1)=F_0$ exactly. Numerical integration of a signed center must
retain its own integration error; (ACM15) specifies exact integrals.

At every independent coefficient face $(W,S;A)$ the coefficient
inclusions and zero entries act componentwise on these same maps.
The original minimum-section difference is a relation $\chi Q$;
successive monic division constructs
$\chi(\sum s_i)Q(\sum s_i)=\sum_i h(s_i)Q_i(\mathbf s)$.
Its original primitive is
$\sum_i(-1)^{i-1}Q_i(D_1,\ldots,D_k)
 (F_h^{\otimes(i-1)}\otimes\phi_*\otimes F_h^{\otimes(k-i)})$.
The tensor differential contributes the same sign and therefore
recovers the relation. At a proper declared support its residual is
the class in
$(d\mathsf C_{\rm full}^{k-1}\cap\mathsf C_{W,S;A}^k)/
d\mathsf C_{W,S;A}^{k-1}$, with the inclusion-induced map and inverse
given by the same representative, as proved in (ISM38)--(ISM43).
Consequently none of the metric estimates deletes this residual,
the full diagonal degree-zero source, or an empty mask's outer label.
